\chapter{High Level Design}
\label{ch:hld}

The high-level design describes the architecture of the benchmark system and the
flow of data through it. This chapter states the design considerations and
development method, the architecture strategy and language choice, the overall
system architecture, and the data flow at successive levels of detail. It closes
with a summary of the design framework.

%----------------------------------------------------------

\section{Design Considerations}

\subsection{General Considerations}

Four considerations shape the design. The first is reproducibility: every
validation decision must be traceable to the inputs that produced it, so each
stage reads and writes schema-versioned artifacts. The second is evidence: a
positive is kept only with a matched sink pattern and a changed sink line. The
third is isolation, which keeps the stages independent, so the corpus, the
mutators, and the harness can be re-run on their own. The fourth is
auditability. The pattern set, the audit reports, and the rejection manifest are
released, which lets a reader check the label decisions instead of trusting
them.

\subsection{Development Method}

The system was built incrementally and validated by audit at each round rather
than once at the end. The three-round audit drove the design: each round
measured the false-positive rate of the labels on disjoint samples, and the
findings of one round motivated the filter added in the next. This
audit-then-tighten loop is the development method, and it is why the false
positive rate fell across rounds rather than being asserted at the start.

\section{Architecture Strategies}

\subsection{Programming Language}

Python is the implementation language for three reasons. The benchmark targets
Python code, so the sink patterns, the AST transformations, and the
semantics-preservation checks all operate in the same language they analyze. The
standard \texttt{ast} module gives a stable, well-documented syntax tree, which
the mutators and the negative builder transform directly. And the machine
learning stack the trained detector needs, PyTorch and Transformers, is native
to Python.

\subsection{Future Plans}

The architecture leaves room for extension. The single-file format can be
replaced by an application-level format that supplies entry points to the
whole-project dataflow detectors, the sink pattern set can be re-derived for
another language, and the mutator suite can gain a dynamic-dispatch axis. Each
extension fits the existing stage boundaries without redesign.

\section{System Architecture}

The system is a pipeline of three stages connected by artifacts. The corpus
stage narrows a raw feed through three validation layers; the mutation stage
generates semantics-preserving variants of the test samples; and the evaluation
stage runs detectors and computes metrics. Figure~\ref{fig:pipeline} shows the
corpus stage, which is the part that establishes label quality. A raw feed of
1{,}982 samples is narrowed by the three layers, and the audited false-positive
rate falls from 52\% to 26\%.

\begin{figure}[htb]
\centering
\resizebox{\textwidth}{!}{%
\begin{tikzpicture}[
  node distance=8mm and 11mm,
  stage/.style={draw, rounded corners, align=center, fill=gray!8,
    minimum height=13mm, text width=23mm, font=\small},
  data/.style={draw, align=center, minimum height=13mm,
    text width=19mm, font=\small},
  arr/.style={-{Stealth[length=2.5mm]}, thick},
  note/.style={font=\footnotesize, align=center}
]
\node[data] (raw) {raw feed\\\textbf{1{,}982}};
\node[stage, right=of raw] (l1) {Layer 1\\sink-presence\\filter};
\node[data, right=of l1] (d1) {\textbf{1{,}233}};
\node[stage, right=of d1] (l2) {Layer 2\\3-round\\adversarial audit};
\node[stage, right=of l2] (l3) {Layer 3\\diff-based\\filter};
\node[data, right=of l3] (final) {\textbf{440} pos.\\$+$\,\textbf{429} safe\\$=$\,\textbf{869}};
\draw[arr] (raw)--(l1);
\draw[arr] (l1)--(d1) node[midway,above,note]{$-749$};
\draw[arr] (d1)--(l2);
\draw[arr] (l2)--(l3);
\draw[arr] (l3)--(final) node[midway,above,note]{$-316$};
\node[note, below=4mm of d1] {audit FP\\\textbf{52\%}};
\node[note, below=4mm of l2] {$\rightarrow$ \textbf{51\%}};
\node[note, below=4mm of l3] {$\rightarrow$ \textbf{26\%}};
\end{tikzpicture}}
\caption{The three-layer label-validation pipeline. Layer~1 rejects samples
lacking the category-defining sink token; Layer~2's independent audit rounds
drive the pattern tightenings; Layer~3 keeps a positive only if the sink line
itself changed between the vulnerable and fixed versions.}
\label{fig:pipeline}
\end{figure}

\section{Data Flow Diagrams}

\subsection{Data Flow Diagram Level 0.0}

At the context level, the system takes two external inputs and produces one
output. The inputs are the advisory and CVE-fix feeds (GHSA, CVEfixes, NVD, OSV,
and the \texttt{vudenc} corpus) and the detector implementations (SAST binaries,
LLM APIs, and the trained model). The output is a set of scored results: per
detector, per CWE, on clean and mutated code. The system is the benchmark
pipeline that transforms the feeds into a validated, mutated corpus and scores
the detectors against it.

\begin{figure}[htb]
\centering
\resizebox{0.85\textwidth}{!}{%
\begin{tikzpicture}[
  font=\small,
  ext/.style={draw, align=center, fill=gray!4,
    minimum height=12mm, text width=30mm},
  proc/.style={draw, rounded corners, align=center, fill=blue!12,
    minimum height=16mm, text width=34mm},
  arr/.style={-{Stealth[length=2.5mm]}, thick, gray!55},
]
\node[ext] (feeds) {Advisory \&\\CVE-fix Feeds\\{\scriptsize GHSA, NVD, OSV}};
\node[proc, right=24mm of feeds] (sys) {\textbf{0.0}\\Benchmark\\Pipeline};
\node[ext, right=24mm of sys] (out) {Scored Results\\{\scriptsize per detector,\\per CWE}};
\node[ext, below=20mm of feeds] (dets) {Detector Suite\\{\scriptsize SAST, Dataflow,\\LLM, GCB}};
\draw[arr] (feeds) -- (sys);
\draw[arr] (dets.east) -- ++(10mm,0) |- (sys.west);
\draw[arr] (sys) -- (out);
\end{tikzpicture}}
\caption{Data Flow Diagram Level~0.0 (context diagram).}
\label{fig:dfd0}
\end{figure}

\subsection{Data Flow Diagram Level 1.0}

One level down, the pipeline is three processes connected by data stores. The
\emph{construction} process reads the raw feeds, applies the three validation
layers, builds the safe twins, and writes the validated corpus to a data store.
The \emph{mutation} process reads the test split and writes seven mutated
variants of each test sample to a variant store. The \emph{evaluation} process
reads both the clean test set and the variant store, runs each detector, and
writes per-sample predictions and aggregate metrics.

\begin{figure}[htb]
\centering
\resizebox{0.95\textwidth}{!}{%
\begin{tikzpicture}[
  font=\small,
  ext/.style={draw, align=center, fill=gray!4,
    minimum height=12mm, text width=24mm},
  proc/.style={draw, rounded corners, align=center, fill=blue!12,
    minimum height=16mm, text width=28mm},
  store/.style={draw, align=center, fill=gray!15,
    minimum height=10mm, text width=24mm},
  arr/.style={-{Stealth[length=2.5mm]}, thick, gray!55},
]
\node[ext] (feeds) {Raw\\Feeds};
\node[proc, right=14mm of feeds] (cons) {\textbf{1.0}\\Corpus\\Construction};
\node[store, right=12mm of cons] (corpus) {Validated\\Corpus};
\node[proc, right=12mm of corpus] (mut) {\textbf{2.0}\\Mutation};
\node[store, right=12mm of mut] (vars) {Variant\\Store};
\node[proc, right=12mm of vars] (eval) {\textbf{3.0}\\Evaluation};
\node[ext, right=14mm of eval] (metrics) {Scored\\Metrics};
\draw[arr] (feeds) -- (cons);
\draw[arr] (cons) -- (corpus);
\draw[arr] (corpus) -- (mut);
\draw[arr] (mut) -- (vars);
\draw[arr] (vars) -- (eval);
\draw[arr] (eval) -- (metrics);
\draw[arr, dashed] (corpus.south) -- ++(0,-12mm) -| (eval.south)
  node[midway, below, font=\scriptsize] {test split};
\end{tikzpicture}}
\caption{Data Flow Diagram Level~1.0: corpus construction, mutation, and evaluation.}
\label{fig:dfd1}
\end{figure}

\subsection{Data Flow Diagram Level 2.0}

The evaluation process expands into the harness shown in
Figure~\ref{fig:harness}. The same 129-sample clean set and eight mutator
variants of the 67 positives flow into seven detectors across four tiers, and
each detector's output flows into the five bounded metrics.

\begin{figure}[htb]
\centering
\resizebox{0.95\textwidth}{!}{%
\begin{tikzpicture}[
  font=\small,
  box/.style={draw, rounded corners, align=center, fill=gray!8,
    minimum height=9mm, text width=34mm},
  src/.style={draw, align=center, fill=gray!4, minimum height=26mm,
    text width=27mm},
  met/.style={draw, align=center, fill=gray!4, minimum height=26mm,
    text width=24mm},
  arr/.style={-{Stealth[length=2.5mm]}, thick, gray!55},
]
\node[src] (bench) {\textbf{Benchmark}\\129 clean\\(67 pos $+$ 62 safe)\\[3pt]$+$ 8 mutator\\variants of the\\67 positives};
\node[box, right=12mm of bench, yshift=17mm] (sast) {\textbf{SAST}\\Bandit, Semgrep};
\node[box, below=2.5mm of sast] (wp) {\textbf{Whole-project}\\Snyk Code, CodeQL};
\node[box, below=2.5mm of wp] (llm) {\textbf{LLM}\\Sonnet 4.6, GPT-4o,\\DeepSeek R1};
\node[box, below=2.5mm of llm] (trained) {\textbf{Trained}\\GraphCodeBERT};
\node[met, right=12mm of llm, yshift=8.5mm] (metrics) {\textbf{Metrics}\\macro-F1\\detection MCC\\SWR\\HCMA\\weighted $\kappa$};
\draw[arr] (bench.east) -- (sast.west);
\draw[arr] (bench.east) -- (wp.west);
\draw[arr] (bench.east) -- (llm.west);
\draw[arr] (bench.east) -- (trained.west);
\draw[arr] (sast.east) -- (metrics.west);
\draw[arr] (wp.east) -- (metrics.west);
\draw[arr] (llm.east) -- (metrics.west);
\draw[arr] (trained.east) -- (metrics.west);
\end{tikzpicture}}
\caption{Data Flow Diagram Level~2.0: the 129-sample benchmark and eight mutator
variants of the 67 positives flow into seven detectors across four tiers,
scored under five bounded metrics.}
\label{fig:harness}
\end{figure}

\section{Summary}

The high-level design is a three-stage, artifact-connected pipeline built around
reproducibility, evidence-backed labels, stage isolation, and auditability. The
system architecture centers on the three-layer validation funnel, and the data
flow refines from a two-input context diagram to the detailed evaluation
harness. Chapter~\ref{ch:detailed} details the design of each module.
